\documentclass{article}
\usepackage{graphicx} % Required for inserting images
\usepackage[russian]{babel}
\usepackage{hyperref}
\title{Мини-конспект по теме: Теорема Пифагора}
\author{Автор: Я Сам }
\date{16 июля 2025 г.}

\begin{document}

\maketitle
\paragraph{Содержание}
\tableofcontents{}

\newpage

\section{Введение}
Теорема Пифагора - одна из важнейших теорем евклидовой геометрии. Она находит применение в самых разных областях:
\begin{itemize}
    \item геометрия
    \item физика
    \item инженерные расчёты
    \item компьютерная графика
\end{itemize}
\section{Формулировка теоремы}
\textbf{Слова:} В прямоугольном треугольнике квадрат гипотенузы равен сумме квадратов катетов.
\begin{center}
    \begin{equation}
    c^2 = a^2 + b^2
    \end{equation}
    \\
    Как видно из формулы 1, знание двух сторон позволяет найти третью
\end{center}
\section{Доказательство (набросок)}
\begin{center}
    Одно из доказательств основывается на площади квадрата, составленного из четырёх одинаковых прямоугольных треугольников и малого квадрата в центре. Раскладывая площадь двумя способами, получаем $c^2 = a^2 + b^2$.
\end{center}
\section{Примеры расчёта}
\paragraph{Пример 1}
\begin{center}
    a = 3, b = 4\\
    $c = \sqrt{a^2 + b^2} = \sqrt{9 + 16} = 5$
\end{center}
\paragraph{Пример 2}
\begin{enumerate}
    \item Дано: a = 5, b = 12
    \item Решение:\\
    \begin{center}
        $c = \sqrt{5^2 + 12^2} = \sqrt{25 + 144} = 13$
    \end{center}
\end{enumerate}
\section{Таблица значений}
\begin{center}
    \begin{tabular}{|c|c|c|}
        \hline
            Катет a&Катет b&Гипотенуза c\\
            \hline
            3&4&5\\
            \hline
            5&12&13\\
            \hline
            7&24&25\\
        \hline
    \end{tabular}
\end{center}
\section{Иллюстрация}
Ниже пример изображения:
\begin{center}
    \includegraphics[]{images.jpg}
\end{center}
\section{Заключение}
Теорема Пифагора - один из краеугольных камней геометрии, помогающий решать множество практических задач.
\section{Ссылки и литература}
\begin{itemize}
    \item Википедия: \href{https://ru.wikipedia.org/wiki/Криштиану_Роналду}{Теорема Пифагора}
    \item Классические учебники геометрии
\end{itemize}
\end{document}
